% TRANSLATION UNIT: translate this file only.
% ===========================================================================
%  11 — Pruebas
%  The \section title is the canonical (Spanish) chapter name; the body prose
%  is written in ENGLISH and translated to Spanish later.
% ===========================================================================

% EVIDENCE INVENTORY (every number below maps to a real artifact):
% EVIDENCE: 311/311 unit tests passing → frontend `npm test` (Karma/ChromeHeadless); CI job "TS — frontend" runs `npm test` (.github/workflows/ci.yml).
% EVIDENCE: coverage 80.3 % statements (1,488/1,853), 64.34 % branches (471/732), 74.42 % functions (422/567), 80.25 % lines (1,427/1,778) → Karma coverage report, frontend, measured 2026-11-03 on the development branch.
% EVIDENCE: 59/59 Playwright specs, 0 skipped → frontend/e2e/** run against the full stack (frontend/playwright.config.ts, testDir ./e2e).
% EVIDENCE: five widths 360/480/768/1024/1280 → WIDTHS array in frontend/e2e/visual/*.spec.ts; breakpoint tokens in frontend/src/styles/tokens/breakpoints.css.
% EVIDENCE: 110 of 112 assertions select by role and accessible name → getByRole/getByLabel count across frontend/e2e, measured 2026-11-03.
% EVIDENCE: pre-redesign baseline 58 passed / 1 conditionally skipped → development-branch e2e run before the redesign.
% EVIDENCE: load profile 200 VUs, 1 min ramp, 3 min hold, 30 s down, 200 pre-seeded accounts → deploy/loadtest/k6-scenarios.js + deploy/loadtest/README.md (VUS=200 RAMP=1m HOLD=3m DOWN=30s).
% EVIDENCE: 64,268 requests at ~259 req/s over 4 m 08 s, 0.00 % HTTP failures (0 of 64,268), 10,478 iterations, read p95 115 ms, action p95 763 ms, zero 4xx/5xx → k6 --summary-export of the attributed run, 2026-11-03.
% EVIDENCE: thresholds read < 1,000 ms / action < 2,000 ms / failures < 1 % → thresholds block in deploy/loadtest/k6-scenarios.js.
% EVIDENCE: Lighthouse 360 px catalog performance 62, LCP 7.4 s; CLS 1280 px 0.383→0.011; JS 77,328→23,022 B gz; font 488,026 B TTF→66,544 B woff2; icon 314,600→240,100 B → Lighthouse runs before/after the fixes, 2026-11-03.
% EVIDENCE: design debt 0 artisanal-class screens, 0 inline styles → node frontend/scripts/design-debt.mjs --check.
% EVIDENCE: prose barrier 0 errors → node docs/informe/scripts/prose-barrier.mjs --check.
% EVIDENCE: 4 px overflow at 360 px on /simuladores/colombia_especifica, 344 px label in a 320 px column → overflow probe in frontend/e2e/zoom-200.spec.ts (expectNothingIsLost) + visual capture.
% EVIDENCE: --text-muted on inset 3.77:1; --fc-ink-3 on accent tint 4.38:1 vs claimed 4.76:1 → surface table in frontend/src/styles/tokens/colors.css.
% EVIDENCE: quiz spec failed on 5 s timeout under parallel load, passed in isolation at 8.4 s/9.4 s; fix was 10 s budget → expect.timeout comment in frontend/playwright.config.ts.
% EVIDENCE: 0.21 % of requests 429 during setup(), one IP vs 600 rpm per-IP limiter → "De dónde salieron los 160 fallos" in deploy/loadtest/README.md.
% EVIDENCE: money-truncation probe reported 0 truncations before/after because clipping was container overflow-x, not text-overflow → expectMoneyIsNeverTruncated in frontend/e2e/visual/simulators.spec.ts.
% EVIDENCE: seed_regression guards FR-049 (engine vs native code oracle) → services/simulator/tests/seed_regression.rs.

\section{Pruebas}

% DONE: verification strategy and measured results, with the defects that only
% an instrument — not human review — could catch, and the rule that a measurement
% artifact must be attributed before it is quoted.

\subsection{Verification strategy}

This chapter reports the verification and technical measurement of FintCart: the instruments that
run against the real stack, the results each of them produced, and the defects that only an
instrument could catch. It is deliberately not the chapter about usability validation. Evaluating
the interface with human participants — task success, the System Usability Scale, and perceived
satisfaction — uses the instruments described in the companion validation package, and its results
belong to chapter~13. The boundary is anchored in the definition of usability itself: ISO~9241-11
decomposes usability into effectiveness, efficiency, and satisfaction
\parencite{iso2018usability}\todo{verify the wording against ISO 9241-11:2018 clause 3.1; the
standard is not in the corpus and this definition is quoted from secondary sources},
and none of those three is measured here, because this chapter does not put a person in front of the
product. What follows is reproducible by a machine: suites, thresholds, and gates.

The verification surface mirrors the architecture, one suite per service plus the cross-cutting
gates. The continuous integration workflow runs, in order: contract validation with \texttt{buf}
(lint plus breaking-change detection); for the Go services, formatting, lint with the prohibition
on floating-point money enforced as a lint rule, and unit tests under the race detector, with the
audit service additionally exercising its immutable, partitioned log against a real PostgreSQL via
testcontainers; for the simulator, formatting, \texttt{clippy} with the \texttt{f32}/\texttt{f64}
prohibition promoted to a hard error, and the full Rust suite, which includes the numeric edge cases
and the seed regression described below; for the TypeScript services and the frontend, lint with the
\texttt{number}-for-money prohibition, the build, and the test suite. Two further jobs verify that
\texttt{dev/build}, \texttt{dev/up}, and \texttt{dev/migrate} still work with zero manual steps, and
that the documented commands and the scripts that implement them have not drifted apart.

One suite deserves its own mention because it guards a decision, not a feature. The simulator's seed
regression suite (\verb|services/simulator/tests/seed_regression.rs|) holds FR-049: the five
calculators of the original requirement are re-seeded on the formula engine while reproducing their
result, so the suite executes the engine and the legacy native code with the same inputs and demands
the same output map. The native modules were not deleted until that suite passed, and it now keeps
them as a living oracle rather than a frozen table. It also defends against its own false green: it
counts how many cases actually reached a computation, because a suite that passes by rejecting every
input is a silent failure wearing a green badge.

\subsection{Unit tests and coverage}

The frontend unit suite runs on Karma in headless Chrome: 311 of 311 tests pass. Coverage is
measured at the same time and reported without a gate, for a reason that has to be stated plainly.
The numbers are recorded in Table~\ref{tab:unit-coverage}.

\begin{table}[htbp]
  \centering
  \begin{tabular}{lrr}
    \hline
    Measure & Covered & Rate \\
    \hline
    Statements & 1,488 / 1,853 & 80.3~\% \\
    Branches    & 471 / 732      & 64.34~\% \\
    Functions   & 422 / 567      & 74.42~\% \\
    Lines       & 1,427 / 1,778  & 80.25~\% \\
    \hline
  \end{tabular}
  \caption{Frontend unit coverage (Karma, ChromeHeadless), development branch, 2026-11-03.}
  \label{tab:unit-coverage}
\end{table}

No coverage threshold is enforced, and the chapter must say why rather than let the gap be
discovered by a reviewer. A build-breaking percentage on this suite would only teach the people
writing it to write filler tests — assertions that touch lines without asserting anything about
them — because that is the cheapest way to move a coverage number. Coverage is therefore read as a
diagnostic, not a contract: the branch figure of 64.34~\% says where behavior is least exercised and
points the next suite at those paths, while the gate that actually fails the build is the set of
behaviors the tests assert. This is a deliberate asymmetry, and it is a limitation, reported as such
in the final section.

\subsection{End-to-end tests}

The end-to-end suite is Playwright driving Chromium against the complete stack raised by
\texttt{dev/up} — gateway, services, and MailHog — not against a frontend-only server, because a
spec that cannot reach a real endpoint fails in a way that looks like a product defect. Fifty-nine
specs pass, zero are skipped. Each visual spec runs its screen at five widths (360, 480, 768, 1024,
and 1280 px), the four breakpoint tokens plus the declared 360 px floor. The numbers are in
Table~\ref{tab:e2e}.

\begin{table}[htbp]
  \centering
  \begin{tabular}{ll}
    \hline
    Aspect & Value \\
    \hline
    Specs & 59 passed, 0 skipped \\
    Widths exercised & 360, 480, 768, 1024, 1280 px \\
    Locator style & role and accessible name (110 of 112 assertions) \\
    Pre-redesign baseline & 58 passed, 1 conditionally skipped \\
    Assertion budget & 10 s (raised from 5 s; see Section on artifacts) \\
    \hline
  \end{tabular}
  \caption{End-to-end suite (Playwright, Chromium) against the full stack.}
  \label{tab:e2e}
\end{table}

The locator style is the point of the suite, not an incidental detail. One hundred and ten of the
112 locator assertions select by role and by accessible name — \verb|getByRole| and
\verb|getByLabel| — rather than by CSS class or DOM position. That is what lets the same assertions
survive a wholesale redesign of the presentation layer, and it is what would make a failure loud:
a test that no longer finds a button by its announced name is reporting an accessibility break, not
a styling change. The two assertions that do not follow the convention are the geometry checks that
only make sense against computed layout, such as the horizontal-overflow probe.

\subsection{Load testing}

The load instrument is a k6 scenario that drives the platform through its real public edge: catalog,
progress, and history reads, and the account-scoped actions, over the OAuth2 authorization-code
flow with PKCE. It runs against 200 pre-seeded accounts — one per virtual user, because the edge
rate-limits per IP and per authenticated user at 600 requests per minute each, and sharing accounts
would measure the limiter instead of the system. The reported run used 200 virtual users, a
one-minute ramp, a three-minute hold, and a thirty-second ramp-down. Its summary is in
Table~\ref{tab:load}.

\begin{table}[htbp]
  \centering
  \begin{tabular}{ll}
    \hline
    Metric & Value \\
    \hline
    Profile & 200 VUs, 1 min ramp, 3 min hold, 30 s down \\
    Accounts & 200 pre-seeded (one per VU) \\
    Duration & 4 m 08 s \\
    Requests & 64,268 (about 259 req/s) \\
    HTTP failures & 0.00~\% (0 of 64,268) \\
    Iterations & 10,478 \\
    Read p95 & 115 ms (threshold: under 1,000 ms) \\
    Action p95 & 763 ms (threshold: under 2,000 ms) \\
    Unexpected responses (any 4xx/5xx) & 0 \\
    \hline
  \end{tabular}
  \caption{Load run (k6) against the development host, 2026-11-03.}
  \label{tab:load}
\end{table}

The thresholds come from the success criteria: reads under 1,000 ms and perceived action response
under 2,000 ms, both at the 95th percentile. The run reports 115 ms and 763 ms respectively, with
zero HTTP failures and zero unexpected responses of any status class. The two caveats that keep this
table honest are stated once here and repeated in the limitations: this is a development host, not
the deployment, so it is not a production capacity figure; and the suite asserts only what its
thresholds assert, never a statistical claim about a population of users.

\subsection{Accessibility, performance, and style gates}

Four gates operate on what a test of behavior cannot see: computed contrast, rendering performance,
and the two forms of debt the redesign exists to retire.

\textbf{Accessibility.} The accessibility suite does not use an external audit library; it measures
in the browser, because contrast is a property of the computed color, not of the token a developer
wrote. It walks each screen with the keyboard, checks that every form control has an associated,
announceable name, and computes WCAG AA contrast on the resolved colors. This is the gate that
caught the token-layer regression described below, and it is intentionally scoped: it is not a
complete WCAG audit, and it does not claim to be one.

\textbf{Performance.} Lighthouse runs against the catalog at 360 px and against the widest layout.
Before the payload fixes, the mobile catalog scored 62 with a largest-contentful-paint of 7.4 s.
The four deltas that follow that score are in Table~\ref{tab:lighthouse}.

\begin{table}[htbp]
  \centering
  \begin{tabular}{lrr}
    \hline
    Metric & Before & After \\
    \hline
    Cumulative layout shift (1280 px) & 0.383 & 0.011 \\
    JavaScript payload (gzipped) & 77,328 B & 23,022 B \\
    Font download & 488,026 B (TTF) & 66,544 B (woff2) \\
    Icon chunk & 314,600 B & 240,100 B \\
    \hline
  \end{tabular}
  \caption{Lighthouse measurements before and after the payload fixes.}
  \label{tab:lighthouse}
\end{table}

The largest single win is the font: serving the same face as woff2 instead of a raw TTF cuts the
download from 488,026 B to 66,544 B, which is the mechanism behind the layout-shift and
largest-contentful-paint improvements, not a coincidence.

\textbf{Style debt.} The design-debt gate (\verb|frontend/scripts/design-debt.mjs|) counts what the
redesign promised to remove rather than what it added: inline styles and screens still referencing
the artisanal CSS classes the old \verb|styles.scss| defined. It reports zero artisanal-class
screens and zero inline styles, and its check mode fails the build on any non-zero value. It
measures exactly one layer of debt — the layer that duplicated the design system's primitives — and
says nothing about the visual quality of what replaced it, which is why visual capture exists as a
separate step.

\textbf{Prose barrier.} One gate does not measure the product at all: the prose barrier
(\verb|docs/informe/scripts/prose-barrier.mjs|) measures this report. It enforces the glossary's
terminology, American spelling, and citation discipline across the chapters, so that three authors
writing in parallel cannot quietly call the simulator a \enquote{calculation engine} or a student a
\enquote{learner}. It reports zero errors on the corpus. Including it here is not vanity: it is the
same idea as the design-debt gate, applied to the document the reader is holding — the quality of
the verification is bounded by the precision of the prose that reports it.

\subsection{What each gate proves, and what it cannot}

Each number above proves one narrow thing and, by construction, nothing else. A green unit suite
proves that 311 tests pass; it does not prove that any page renders. The project has a concrete
witness for that distinction: a logo mask once passed the build and broke the development server,
because the build compiled the template and never painted it. The build is a necessary gate and a
false comfort, and visual capture — actually reading a rendered screen — is the only gate that
closes that particular hole. Likewise, coverage is a diagnostic without a threshold for the reason
already given, and the load figures are a measurement of this development host on this day, not a
capacity statement about the deployment.

The accessibility gate proves contrast on the surfaces it visits. The performance gate proves the
payload shrank. The style and prose gates prove that a specific kind of debt is zero. None of them
proves that the interface is good, or that anyone learned anything; those claims belong to the
instruments of chapter~13, and this chapter declines them on purpose.

\subsection{Defects that human review missed}

The strongest result in this chapter is negative evidence: two defects shipped past human
inspection and past the ad hoc screenshot sweep, and were caught only by an instrument that asserts
what nobody was looking at. They are the intellectual contribution of the chapter, so they are
reported in the same detail as the green numbers.

\textbf{A four-pixel overflow.} At 360 px, the calculator screen
\verb|/simuladores/colombia_especifica| overflowed horizontally by 4 px: a sentence-length button
label with \verb|white-space: nowrap| measured 344 px inside a 320 px column. The ad hoc capture
never visited that route — the sweep photographed the routes it thought representative, and this one
was not among them — and a 4 px overflow is exactly the defect a person scrolling a screenshot does
not notice. The overflow probe, which compares the document's scroll width against its client width
and fails on a one-pixel difference, noticed immediately.

\textbf{A contrast regression in the token layer.} The accessibility gate measured
\verb|--text-muted| on the inset surface at 3.77:1 — below the WCAG AA minimum of 4.5:1 — and,
later, \verb|--fc-ink-3| on the accent tint at 4.38:1, while the token file claimed 4.76:1. The
claim was not a lie and not an error of arithmetic: it had been measured against white, and the
token is used on a tinted surface. The comment in the token file now records the four surfaces and
their ratios, and the fix darkened the ink rather than lightening the tint, because lightening the
tint would have broken the token on every surface darker than paper. What matters for this chapter
is the failure mode: a number in a source file that was true in one context and false in the one
that shipped, and that only computed-color measurement — not reading the token, not reviewing the
diff — could contradict.

Table~\ref{tab:defects} summarizes both, plus the two near-misses of the next subsection, which are
a different class of finding.

\begin{table}[htbp]
  \centering
  \begin{tabular}{p{3.6cm}p{4.6cm}p{5.6cm}}
    \hline
    Finding & What review saw & What the instrument measured \\
    \hline
    4 px overflow at 360 px & Nothing: the route was not in the ad hoc sweep & A 344 px label inside a 320 px column on \verb|/simuladores/colombia_especifica| \\
    Contrast regression & A token file claiming 4.76:1 & 3.77:1 on the inset surface; 4.38:1 on the accent tint — measured against the surface in use, not against white \\
    Quiz timeout & A product failure & A 5 s assertion timeout under parallel load; passed in isolation at 8.4 s and 9.4 s \\
    429s under load & A rate-limiting failure of the platform & 0.21 \% of requests returning 429 during \verb|setup()|, before the measured window, from one shared IP against the 600 rpm per-IP limiter \\
    \hline
  \end{tabular}
  \caption{Defects and near-misses, and the difference between them.}
  \label{tab:defects}
\end{table}

\subsection{Measurement artifacts are part of the result}

The third class is methodological, and it produces the rule this chapter wants the reader to keep.
Two failures were not product defects at all.

A quiz spec failed on a 5 s assertion timeout under parallel load and passed twice in isolation, at
8.4 s and 9.4 s. The fix was ten seconds of assertion budget, not a line of product code, because
an assertion that fails under load does not measure what it claims to measure: it measures how
quickly the machine answered that day, and a 5 s window was simply too tight for a full-stack spec
competing for the gateway. And 0.21 \% of the requests in one load run returned 429 — every one of
them during \verb|setup()|, before the measured window, because the harness shared a single IP
against a 600 requests-per-minute per-IP limiter. That was a harness bug that would have been
published as a product failure: the number existed, had decimals, and pointed at the wrong thing.

The rule, stated once: \emph{a measurement artifact is part of the result, and must be attributed
before it is quoted.} The load instrument now tags every unexpected response with its phase —
\verb|vu| for the measured window, \verb|setup| for preparation — precisely so that a setup
collision can never again be summed into a headline failure rate. The lesson generalizes: a gate
that cannot tell the difference between a defect in the product and a defect in the instrument that
measures it will eventually report the instrument's defect as the product's.

\subsection{Limitations}

The honest boundaries, in one place. Coverage has no enforced threshold, by design, and therefore
cannot be cited as a guarantee of exercised behavior. The money-truncation probe reported zero
truncations both before and after the clipping was fixed, because the clipping was a container's
\verb|overflow-x|, not \verb|text-overflow| — the probe measured the wrong layer, and only reading
the screenshot found the real defect; a monetary figure rendered as \verb|$1.234.567,89| must be
read whole, because a truncated figure is a false datum, and the probe as written could not see that
violation. The load run happened on a single development host, not on the deployment, so it is not a
production capacity figure and no statistical validity is claimed for it. And no green build in
this chapter proves that a page renders; the logo mask is the standing counterexample.

Taken together, the chapter's claim is modest and specific: the platform carries a verification
surface broad enough that a four-pixel overflow and a false contrast claim could not survive it,
and the same surface, when it failed, failed informatively enough to distinguish a product defect
from a measurement artifact. That is what testing is for, and that is all this chapter claims it
achieved.
